\subsection{Multilinearization and the exact kernel-fold functional}\label{sec:binary:multilinearization}
We retain $\ell\geq1$, $n=2^\ell$, and a characteristic-two field $\F$.
For $0\leq e<n$ write its unique binary expansion as
$e=\sum_{i=0}^{\ell-1}e_i2^i$ with $e_i\in\{0,1\}$.

\begin{definition}[Binary-exponent multilinearization]\label{def:binary:ml}
Let $\mathscr{ML}_\ell$ be the $\F$-vector space of
polynomials multilinear in $Z_0,\ldots,Z_{\ell-1}$.
Define the $\F$-linear map
\[
 \mathscr M_\ell:\F[X]_{<n}\longrightarrow\mathscr{ML}_\ell,
 \qquad
 \mathscr M_\ell\Bigl(\sum_{e=0}^{n-1}a_eX^e\Bigr)
 :=\sum_{e=0}^{n-1}a_e\prod_{i=0}^{\ell-1}Z_i^{e_i}.
\]
\end{definition}

\begin{lemma}[Multilinearization is a vector-space isomorphism]\label{lem:binary:ml-iso}
The map $\mathscr M_\ell$ is an $\F$-linear bijection.
\end{lemma}
\begin{proof}
Linearity is part of its definition. The monomials
$X^e$, $0\leq e<n$, form a basis of $\F[X]_{<n}$.
The monomials $\prod_iZ_i^{e_i}$ indexed by the binary digits of
$0\leq e<n$ form a basis of $\mathscr{ML}_\ell$.
The binary expansion gives a bijection between these two bases.
The map sends each element of the first basis to the corresponding
element of the second, hence is an isomorphism.
\end{proof}

\begin{lemma}[Explicit subset expansion of a kernel]\label{lem:binary:kernel-expansion}
For $b\in I_\ell$,
\[
 K_b(X)=\sum_{S\subseteq\{0,\ldots,\ell-1\}}
   \left(\prod_{i\notin S}b_i\right)
   X^{\sum_{i\in S}2^i}.
\]
\end{lemma}
\begin{proof}
In the product $K_b=\prod_i(X^{2^i}+b_i)$,
a term is obtained by selecting $X^{2^i}$ from a subset $S$
of factors and $b_i$ from the complementary factors.
The resulting term has scalar coefficient $\prod_{i\notin S}b_i$
and monomial factor $X^{\sum_{i\in S}2^i}$.
Every subset occurs exactly once in the full distributive
expansion, establishing the equality.
\end{proof}

\begin{lemma}[Multilinearization of every kernel basis element]\label{lem:binary:ml-kernel}
For every $b\in I_\ell$,
\[
 \boxed{\mathscr M_\ell(K_b)(Z)=\prod_{i=0}^{\ell-1}(Z_i+b_i).}
\]
\end{lemma}
\begin{proof}
By Lemma~\ref{lem:binary:kernel-expansion} and linearity,
\[
 \mathscr M_\ell(K_b)=\sum_{S}
   \left(\prod_{i\notin S}b_i\right)
   \mathscr M_\ell\left(X^{\sum_{i\in S}2^i}\right).
\]
The binary digits of $\sum_{i\in S}2^i$ equal $1$ precisely for
$i\in S$, so the last factor is $\prod_{i\in S}Z_i$.
Therefore the sum equals
$\sum_S(\prod_{i\notin S}b_i)(\prod_{i\in S}Z_i)$.
The distributive expansion of $\prod_i(Z_i+b_i)$ is the same sum.
This argument uses only linearity of $\mathscr M_\ell$,
not multiplicativity.
\end{proof}

\begin{definition}[Characteristic-two equality weights]\label{def:binary:lambda}
For $b\in I_\ell$ and $\theta=(\theta_0,\ldots,\theta_{\ell-1})\in\F^\ell$,
define
\[
 \lambda_b(\theta):=\prod_{i=0}^{\ell-1}(1+\theta_i+b_i).
\]
For $P=\sum_bc_bK_b$ with its unique kernel coefficient table,
define
\[
 D_P(\theta):=\sum_bc_b\lambda_b(\theta).
\]
\end{definition}

\begin{lemma}[Equality weights agree]\label{lem:binary:eq-weights}
For all $b\in I_\ell$ and $\theta\in\F^\ell$,
$\lambda_b(\theta)=\operatorname{eq}_\ell(b,\theta)$.
\end{lemma}
\begin{proof}
For each coordinate $i$, if $b_i=0$ then
$1+\theta_i+b_i=1+\theta_i=1-\theta_i$.
If $b_i=1$ then $1+\theta_i+b_i=\theta_i$,
because $1+1=0$. These are precisely the two cases of
$(1-b_i)(1-\theta_i)+b_i\theta_i$.
Multiply the coordinatewise identities.
\end{proof}

\begin{lemma}[Boolean weights partition unity]\label{lem:binary:lambda-partition}
For every $\theta\in\F^\ell$,
$\sum_{b\in I_\ell}\lambda_b(\theta)=1$.
\end{lemma}
\begin{proof}
Factor the sum over bit strings into a product of two-term sums:
\[
 \sum_b\lambda_b(\theta)
 =\prod_{i=0}^{\ell-1}
  \bigl((1+\theta_i+0)+(1+\theta_i+1)\bigr)
 =\prod_i\bigl((1+\theta_i)+\theta_i\bigr)=1.
\]
\end{proof}

\begin{theorem}[Complete kernel fold equals multilinearized evaluation]\label{thm:binary:fold-ml}
For every $P\in\F[X]_{<n}$ and $\theta\in\F^\ell$,
\[
 \boxed{
 \mathsf{PFold}^{[\ell]}_\theta(P)
  =D_P(\theta)
  =\mathscr M_\ell(P)(1+\theta_0,\ldots,1+\theta_{\ell-1}).
 }
\]
\end{theorem}
\begin{proof}
By Proposition~\ref{prop:intro:basis}, uniquely write
$P=\sum_bc_bK_b$. The fold dictionary,
Corollary~\ref{cor:kbff:full-evaluation},
identifies its complete fold with
$\sum_bc_b\operatorname{eq}_\ell(b,\theta)$,
which equals $D_P(\theta)$ by Lemma~\ref{lem:binary:eq-weights}.
For the second equality, apply the linear map $\mathscr M_\ell$
to the kernel expansion of $P$ and evaluate:
\[
 \mathscr M_\ell(P)(1+\theta)
 =\sum_bc_b\,\mathscr M_\ell(K_b)(1+\theta).
\]
Lemma~\ref{lem:binary:ml-kernel} evaluates each summand as
$\prod_i(1+\theta_i+b_i)=\lambda_b(\theta)$.
Thus the sum is precisely $D_P(\theta)$.
\end{proof}

\begin{definition}[Partial multilinearization]\label{def:binary:partial-ml}
For $0\leq j\leq\ell$, define
\[
 \mathscr M_{\ell,j}(P)(Z_0,\ldots,Z_{j-1},Y)
 :=\sum_{e=0}^{n-1}a_e
   \left(\prod_{i=0}^{j-1}Z_i^{e_i}\right)
   Y^{\lfloor e/2^j\rfloor}
\]
for $P(X)=\sum_{e=0}^{n-1}a_eX^e$.
\end{definition}

\begin{lemma}[Partial folding formula]\label{lem:binary:partial-fold}
For $0\leq j\leq\ell$ and $r\in\F^j$,
\[
 \boxed{
 \mathsf{PFold}^{[j]}_r(P)(Y)
 =\mathscr M_{\ell,j}(P)(1+r_0,\ldots,1+r_{j-1},Y).
 }
\]
\end{lemma}
\begin{proof}
The explicit fold formula of
Theorem~\ref{thm:kbff:fold-coefficients}
gives a summand
$a_e\bigl(\prod_{i<j}A_{e_i}(r_i)\bigr)Y^{\lfloor e/2^j\rfloor}$.
In characteristic two, $A_0(r_i)=2r_i-1=1$ and
$A_1(r_i)=1-r_i=1+r_i$.
Thus $\prod_{i<j}A_{e_i}(r_i)=\prod_{i<j}(1+r_i)^{e_i}$.
Substitution produces exactly the evaluation of
Definition~\ref{def:binary:partial-ml}.
\end{proof}

\begin{lemma}[Frobenius covariance of equality weights]\label{lem:binary:lambda-frob}
For $b\in I_\ell$, $\theta\in\F^\ell$, and $j\geq0$,
\[
 \lambda_b(\theta)^{2^j}
 =\lambda_b(\theta_0^{2^j},\ldots,\theta_{\ell-1}^{2^j}).
\]
\end{lemma}
\begin{proof}
Apply $\sigma_j$ to the product
$\prod_i(1+\theta_i+b_i)$, using the ring-homomorphism
property of Lemma~\ref{lem:binary:frob-ring}.
Since $1$ and $b_i\in\F_2$ are fixed by $\sigma_j$,
each factor becomes $1+\theta_i^{2^j}+b_i$.
Their product is the equality weight on the right.
\end{proof}

\begin{corollary}[Frobenius covariance of kernel-fold values]\label{cor:binary:fold-frob}
For $c:I_\ell\to\F$, $\theta\in\F^\ell$, and $j\geq0$,
\[
 \boxed{D_{U_c}(\theta)^{2^j}
 =D_{U_{c^{\sigma_j}}}(\theta^{\sigma_j}).}
\]
\end{corollary}
\begin{proof}
Raise the finite sum
$D_{U_c}(\theta)=\sum_bc(b)\lambda_b(\theta)$
to its $2^j$-th power.
By Lemma~\ref{lem:binary:frob-add} and multiplicativity,
it becomes $\sum_bc(b)^{2^j}\lambda_b(\theta)^{2^j}$.
Apply Lemma~\ref{lem:binary:lambda-frob}
and Definition~\ref{def:binary:table-frob}.
\end{proof}

\begin{remark}[Nonmultiplicativity warning]
The multilinearization map is \emph{not} a ring homomorphism,
even where both sides happen to be defined.
For $\ell\geq2$,
$\mathscr M_\ell(X^2)=Z_1$ while
$\mathscr M_\ell(X)\mathscr M_\ell(X)=Z_0^2$.
The polynomials $Z_1$ and $Z_0^2$ are distinct.
All foregoing kernel identities were proved from expansions
and the linearity of $\mathscr M_\ell$.
\end{remark}
